% !TEX program = lualatex
\documentclass[11pt,a4paper]{article}

\usepackage{luatexja}
\usepackage{amsmath,amssymb,amsthm}
\usepackage{geometry}
\geometry{margin=1in}

%-------------------------------------------------
%  定理環境
%-------------------------------------------------
\theoremstyle{definition}
\newtheorem{definition}{定義}[section]
\newtheorem{axiom}{公理}[section]
\newtheorem{theorem}{定理}[section]
\newtheorem{corollary}{系}[section]
\newtheorem{principle}{原理}[section]

%-------------------------------------------------
%  独自マクロ
%-------------------------------------------------
\newcommand{\dfix}{D_{\mathrm{fix}0}} % 0次元不動点（空）
\newcommand{\mweq}{\mathrel{\overset{\mathrm{mw}}{=}}}
\newcommand{\metanegation}{\Uparrow}
\newcommand{\susp}[1]{#1^*}
\newcommand{\aletheiai}{\Vdash} % 顕現関係
\newcommand{\Leadsto}{\rightsquigarrow} % 状態遷移
\newcommand{\Cat}{\mathbf{C}_{\mathrm{cat}}}  % 圏論的空間
\newcommand{\Bou}{\mathbf{B}_{\mathrm{bou}}}  % 界論的空間

\renewcommand{\abstractname}{Abstract}

\begin{document}

\title{帰納フラクタロイドの定式化：\\
論理体系から実行可能プロセスへの着火機構}
\author{Masaomi Chiba}
\date{2026-04-25}
\maketitle

\begin{abstract}
本稿では、任意の静的な論理・文法体系 $N$（世俗諦）を、
界論（Bou）上の動的な実行プロセス（勝義諦）へと変換する普遍的機構として
「帰納フラクタロイド」$I_N$ を導入する。
従来の計算機科学における「モナド」は自己同一性（$id$）を要求するため
本体系には適合しない。界論における $I_N$ は、論理の「隙間」を埋めるのではなく、
そこに内在する「無明（Avidy\=a）」を動力源として顕現の連鎖（$\aletheiai$）を
駆動し、最終的に無記状態 $\dfix$ へと還滅させるフラクタルな変換器である。
\end{abstract}

\section{帰納フラクタロイドの定義}

あらゆる形式体系には「記述」と「実行」の間に固有の隙間（無明の温床）が存在する。
これを界論の動的プロセスへ接続するための機構を定義する。

\begin{definition}[帰納フラクタロイド $I_N$]
論理または文法体系 $N$ に対し、帰納フラクタロイド $I_N$ とは、
$N$ を界 $\Bou$ 上の動的実行プロセスへと変換する文脈演算子であり、
以下の機能を持つ。
\begin{enumerate}
    \item \textbf{保留（Suspension）}: 静的対象 $x \in N$ を、界論の実行コンテキストへ $\susp{x}$ として埋め込む（無明の発生）。
    \item \textbf{展開（Manifestation）}: $N$ に欠落していた制御構造（再帰・状態・探索など）を、メタ否定 $\metanegation$ と顕現関係 $\aletheiai$ の連鎖として展開する。
    \item \textbf{還滅（Resolution）}: 蓄積された無明を消費し、無記状態 $\dfix$ への相転移（大域的カット消去）を導く。
\end{enumerate}
\end{definition}

\begin{definition}[プログラムの結合と着火]
体系 $N$ を帰納フラクタロイド $I_N$ によって実行可能なプロセス
（プログラム）へと着火する操作を $\otimes$ で表す。
$$ \text{Lang}(N) := N \otimes I_N $$
このとき、プログラムの実行は以下の顕現の連鎖として記述される。
$$ \susp{P_0} \aletheiai P_1 \aletheiai \dots \aletheiai \metanegation^4 P_k \mweq \dfix $$
\end{definition}

\section{無明の消費と停止性}

従来のモナド理論では、ノルムや型の制約によって停止性を外部から保証していたが、
帰納フラクタロイドにおいては、無明（自己参照の閉じ損ね）の力学そのものが停止を導く。

\begin{theorem}[無明の還滅による停止性]
実行可能なプログラム $\text{Lang}(N)$ のプロセスにおいて、
帰納フラクタロイド $I_N$ が駆動する各ステップ（$\aletheiai$）は、
系に内在する「構造的無明（Avidy\=a）」を必然的に消費していく。
無明が尽きた時点で四段階収束則が発動し、プログラムは無記状態 $\dfix$
（解脱・計算の完了）へと到達して停止する。
\end{theorem}
\begin{proof}
界論の公理に基づく中道推論規則により、すべての顕現プロセスは $\dfix$ を共通の
到達点（Sink）として持つ。$I_N$ はこの自然な「沈み込み（重力）」に沿って
状態を遷移させるため、無明の量は有限ステップでゼロに収束し、$P_k \mweq \dfix$ が成立する。
\end{proof}

\section{4つの典型例の再解釈}

旧来の「モナドによる補完」は、界論においては「無明の形態に応じたフラクタル展開」
として記述される。

\begin{corollary}[ペアノ算術と再帰フラクタロイド]
ペアノ算術 $\mathbb{P}$ に欠落している「計算の合成」という無明は、
再帰構造を展開するフラクタロイド $I_{\mathbb{P}}$（旧・自由モナド）によって駆動され、
$\mathbb{P} \otimes I_{\mathbb{P}}$ は動的な\textbf{項書換え系}となる。
\end{corollary}

\begin{corollary}[最小論理と継続フラクタロイド]
最小論理 $\mathrm{ML}$ に欠落している「脱出・否定」という無明は、
計算のコンテキスト自体をメタ化（$\metanegation$）して保持する
継続フラクタロイド $I_{\mathrm{ML}}$ によって駆動され、
$\mathrm{ML} \otimes I_{\mathrm{ML}}$ は call/cc を備えた\textbf{古典論理の実行系}となる。
\end{corollary}

\begin{corollary}[文脈自由文法と状態フラクタロイド]
文脈自由文法 $\mathrm{CFG}$ に欠落している「位置の記憶」という無明は、
状態を暗黙の引数として連鎖（$\aletheiai$）させる
状態フラクタロイド $I_{\mathrm{CFG}}$ によって駆動され、
$\mathrm{CFG} \otimes I_{\mathrm{CFG}}$ は\textbf{再帰下降パーザ}として駆動する。
\end{corollary}

\begin{corollary}[ホーン節と探索フラクタロイド]
ホーン節 $\mathrm{Horn}$ に欠落している「選択とバックトラック」という無明は、
非決定的な分岐を並行して顕現させる
探索フラクタロイド $I_{\mathrm{Horn}}$（旧・リストモナド）によって駆動され、
$\mathrm{Horn} \otimes I_{\mathrm{Horn}}$ は\textbf{Prolog実行系}となる。
\end{corollary}

\section*{結論}

「帰納フラクタロイド」の概念は、論理体系とプログラミング言語の間の隙間が、
バグではなく「無明（計算を駆動するためのポテンシャルエネルギー）」であること
を明確にする。自己同一性（モナド則）の呪縛から解放された本体系は、
あらゆる論理・文法を、$\dfix$ に至るフラクタルな「顕現と還滅の力学系」
として統一的にコンパイルすることを可能にする。

\end{document}
